% संपूर्ण मराठी ग्रंथातील प्रकरण 30: उपपत्ती आणि संगणनक्षमता
% हा संपादनयोग्य स्रोतखंड आहे. संकलनासाठी संपूर्ण स्रोतसंग्रहातील
% अक्षररूपे, पूर्वभाग, चिन्हव्याख्या आणि आकृती-स्रोत आवश्यक आहेत.
% मूळ: Open Logic Project; मजकूर आणि रूपांतर CC BY 4.0.
% यंत्रानुवाद, दुरुस्ती आणि तपासणी: OpenAI Codex —
% GPT-5.6 Sol आणि GPT-6 Sol, दोन्ही Ultra effort.
% दुरुस्ती-पुनरावलोकन: GPT-6.1 Sol, Ultra effort.
\chapter{उपपत्ती आणि संगणनक्षमता}\label{inc:tcp:chap}
\begin{editorial}
  हे प्रकरण संगणनक्षमता सिद्धांतावरील प्रकरणातील सामग्रीवर
  अवलंबून आहे; ते शिकवले नसेल, तर हे प्रकरण वगळता येईल.
  सध्या हा जेरेमी अविगॅड यांच्या टिपणांवरून केलेला प्राथमिक
  रूपांतरित मसुदा आहे. त्याचे अद्याप पुनरावलोकन झालेले नाही
  आणि सराव-प्रश्नही दिलेले नाहीत.
\end{editorial}









\section{परिचय}\label{inc:tcp:int:sec}

\begin{editorial}
हा विभाग पुन्हा लिहिणे आवश्यक आहे.
\end{editorial}

आपल्याकडे पुढील गोष्टी उपलब्ध आहेत:
\begin{enumerate}
\item एखादे फलन $\Th{Q}$ मध्ये निरूपणीय असणे म्हणजे काय,
  याची व्याख्या (\ref{inc:req:int:defn:representable-fn})
\item एखादा संबंध $\Th{Q}$ मध्ये निरूपणीय असणे म्हणजे काय,
  याची व्याख्या (\ref{inc:req:rel:defn:representing-relations})
\item $\Th{Q}$ मध्ये निरूपणीय फलने नेमकी संगणनक्षम
  फलनेच असतात, असे सांगणारे प्रमेय
  (\ref{inc:req:int:thm:representable-iff-comp})
\item $\Th{Q}$ मध्ये निरूपणीय संबंध नेमके संगणनक्षम
  संबंधच असतात, असे सांगणारे प्रमेय
  (\ref{inc:req:rel:thm:representing-rels})
\end{enumerate}

{ \em उपपत्ती} म्हणजे तार्किक निष्पन्नतेखाली बंद असलेला वाक्यसंच.
म्हणजे $T$ ने $\varphi$ सिद्ध केले, तर $\varphi$ हे $T$ मध्ये
असते. उपपत्तीचा विचार वाक्यांचा आणि त्यांच्यापासून
निष्पन्न होणाऱ्या सर्व वाक्यांचा संग्रह म्हणून करणे
उपयुक्त ठरेल. यापुढे $\Th{Q}$ हे चिन्ह
\ref{inc:req:int:sec} मधील आठ स्वयंसिद्धकांपासून
निष्पन्न होणाऱ्या सर्व वाक्यांच्या { \em उपपत्ती}साठी
वापरू. $\Th{Q}$ मधील सूत्रांना संख्यांनी
सांकेतिक करता येते, हे लक्षात ठेवा. $\varphi$ असे
सूत्र असेल, तर $\varphi$ ला सांकेतिक करणारी
संख्या $\Gn{\varphi}$ ने दर्शवू. या सांकेतीकरणाच्या
आधारे आता सूत्रांचे विविध संच संगणनक्षम आहेत की
नाहीत, हा प्रश्न विचारता येतो.












\section[Q हा संगणनक्षमपणे प्रगणनीय-संपूर्ण आहे]{$\Th{Q}$ हा संगणनक्षमपणे प्रगणनीय-संपूर्ण आहे}\label{inc:tcp:qce:sec}


\begin{thm}
$\Th{Q}$ हा संगणनक्षमपणे प्रगणनीय आहे, पण निर्णेय नाही.
खरे तर तो संपूर्ण संगणनक्षमपणे प्रगणनीय संच आहे.
\end{thm}

\begin{proof}
$\Th{Q}$ हा संगणनक्षमपणे प्रगणनीय आहे हे पाहणे अवघड नाही:
तो अशा वाक्यांच्या संकेतांकांचा संच आहे की
$y$ हा वाक्याचा संकेतांक असेल, तर $\Th{Q}$ मध्ये
$y$ ने सांकेतिक केलेल्या वाक्याची एखादी सिद्धता $x$ अस्तित्वात असते:
\[Q = \Setabs{y}{\lexists[x][\Prf[\Th{Q}](x,y)]}.
\]
पण $\Prf[\Th{Q}](x,y)$ संगणनक्षम आहे
(खरे तर आदिम पुनरावर्ती आहे), आणि वरच्या रूपात
लिहिता येणारा कोणताही संच संगणनक्षमपणे प्रगणनीय असतो.

तो संपूर्ण संगणनक्षमपणे प्रगणनीय संच आहे, असे म्हणणे
$K \leq_m Q$ असे म्हणण्यास समतुल्य आहे. येथे
$K = \Setabs{x}{\cfind{x}(x) \fdefined}$.
म्हणून $K$ चे~$\Th{Q}$ कडे न्यूनीकरण दाखवू.
क्लिनीचे विधेय $T(e,x,s)$ आदिम पुनरावर्ती असल्यामुळे
त्याचे $\Th{Q}$ मध्ये, समजा $\varphi_T$ या सूत्राने, निरूपण
होते. मग प्रत्येक $x$ साठी
\begin{align*}
x \in K & \lif \lexists[s][T(x,x,s)] \\
& \lif \lexists[s][(\Th{Q} \Proves \varphi_T(\num x,\num x, \num s))]
  \\
& \lif \Th{Q} \Proves \lexists[s][\varphi_T(\num x, \num x, s)].
\end{align*}
उलट $\Th{Q} \Proves \lexists[s][\varphi_T(\num x, \num x, s)]$
असेल, तर या सिद्धांताची स्वयंसिद्धके मानक नैसर्गिक-संख्या
प्रतिमानात सत्य असल्याने त्या प्रतिमानात हे अस्तित्ववाचक
वाक्य सत्य असते. म्हणून काही नैसर्गिक संख्या~$n$ साठी
$\varphi_T(\num x, \num x, \num n)$ सत्य असते.
$T(x,x,n)$ असत्य असेल, तर $\varphi_T$ हे $T$ चे
निरूपण करत असल्याने $\Th{Q}$ ही
$\lnot \varphi_T(\num x, \num x, \num n)$ सिद्ध करेल.
मात्र ते मानक प्रतिमानात असत्य सूत्र ठरेल आणि
$\Th{Q}$ ते सिद्ध करेल, हा विरोधाभास आहे.
म्हणून $T(x,x,n)$ सत्य असले पाहिजे; त्यावरून
$\cfind{x}(x) \fdefined$.

थोडक्यात, प्रत्येक $x$ साठी $x$ हा $K$ मध्ये
असतो तर आणि तरच $\Th{Q}$ हे
$\lexists[s][\varphi_T(\num x,\num x,s)]$ सिद्ध करते.
म्हणून $x$ पासून
$\lexists[s][\varphi_T(\num x,
  \num x, s)]$ या वाक्याचा संकेतांक देणारे
$f$ हे फलन $K$ चे~$\Th{Q}$ कडे न्यूनीकरण करते.
संख्यांकाचे प्रतिस्थापन आणि सूत्राचे सांकेतीकरण
संगणनक्षम असल्याने हे न्यूनीकरण संगणनक्षम आहे.
विकर्ण थांबण्याचा संच आधीच अनिर्णेय असल्याने
हे न्यूनीकरण या संचाचाही अनिर्णेयपणा दाखवते.
\end{proof}












\section[Q चा ओमेगा-सुसंगत विस्तार अनिर्णेय असतो]{$\Th{Q}$ चा $\omega$-सुसंगत विस्तार अनिर्णेय असतो}\label{inc:tcp:oqn:sec}

\begin{explain}
$\Th{Q}$ हा c.e.-संपूर्ण आहे या सिद्धतेत,
$\Th{Q}$ मध्ये सिद्ध होणारे प्रत्येक वाक्य नैसर्गिक संख्यांबाबत
``सत्य'' असते या तथ्याचा आधार घेतला होता. पुढील व्याख्या आणि
प्रमेय, वरील सिद्धतेत आवश्यक असलेले ``सत्य''चे नेमके पैलू
ओळखून, त्या प्रमेयाचे सबलीकरण करतात. मात्र या प्रमेयावर फार
रेंगाळू नका, कारण लवकरच आपण त्याहून अधिक सबल प्रमेय सिद्ध करू.
ते मुख्यतः ऐतिहासिक कारणासाठी समाविष्ट केले आहे: गोडेल यांच्या
मूळ लेखात $\omega$-सुसंगततेची संकल्पना वापरली होती; पण लगेचच
$\omega$-सुसंगततेच्या जागी साधी सुसंगतता घालून त्यांचा निकाल
अधिक सबल करण्यात आला.
\end{explain}

\begin{defn}
\label{inc:tcp:oqn:thm:oconsis-q}
उपपत्ती $\Th{T}$ ही $\omega$-सुसंगत असते, जर पुढील अट पूर्ण होत असेल:
$\lexists[x][\varphi(x)]$ हे कोणतेही वाक्य असो. $\Th{T}$ कडून
$\lnot \varphi(\num 0)$, $\lnot \varphi(\num 1)$,
$\lnot \varphi(\num 2)$, \dots ही सर्व वाक्ये सिद्ध होत असतील,
तर $\Th{T}$ कडून $\lexists[x][\varphi(x)]$ सिद्ध होत नाही.
\end{defn}

\begin{thm}
\label{inc:tcp:oqn:thm:oconsis-q-undec}
  $\Th{T}$ ही कोणतीही $\omega$-सुसंगत उपपत्ती असेल आणि तिच्यात
  $\Th{Q}$ समाविष्ट असेल, तर $\Th{T}$ अनिर्णेय असते.
\end{thm}

\begin{proof}
$\Th{T}$ मध्ये $\Th{Q}$ समाविष्ट असल्यास संगणनक्षम
फलनांचे आणि संबंधांचे $\Th{T}$ मध्ये निरूपण होते.
आधीची सिद्धता किंचित बदलली तरी पुरेसे आहे. पूर्वीप्रमाणे,
$x \in K$ असेल, तर $\Th{T}$ हे $\lexists[s][\varphi_T(\num
  x,\num x, s)]$ सिद्ध करते. उलट, $\Th{T}$ हे
$\lexists[s][\varphi_T(\num x, \num x, s)]$ सिद्ध करते असे धरू.
मग $x$ हा $K$ मध्ये असलाच पाहिजे: अन्यथा यंत्र $x$ चे
$x$ या आदानावरील थांबणारे संगणन अस्तित्वात नसते. $\varphi_T$ हे
क्लिनीच्या $T$ संबंधाचे निरूपण करत असल्याने $\Th{T}$ कडून
$\lnot \varphi_T(\num x, \num x, \num 0)$,
$\lnot \varphi_T(\num x, \num x, \num 1)$, \dots
सिद्ध होतील. यामुळे $\Th{T}$ ही $\omega$-विसंगत ठरेल.
\end{proof}












\section[Q चे सुसंगत विस्तार अनिर्णेय असतात]{$\Th{Q}$ चे सुसंगत विस्तार अनिर्णेय असतात}\label{inc:tcp:cqn:sec}

\begin{explain}
कोणत्याही सूत्र~$\varphi$ साठी $\varphi$ आणि $\lnot \varphi$ ही दोन्ही सूत्रे
सिद्ध होत नसतील, तर उपपत्ती \emph{सुसंगत} असते, हे लक्षात ठेवा.
व्याघातापासून कोणतेही वाक्य निष्पन्न होत असल्यामुळे विसंगत उपपत्ती
क्षुल्लक असते: तिच्यात प्रत्येक वाक्य सिद्ध होते.
उपपत्ती $\omega$-सुसंगत असेल, तर ती सुसंगतही असते, हे स्पष्ट आहे.
पण सुसंगतता ही अधिक दुर्बल अट आहे (म्हणजे सुसंगत असलेल्या, पण
$\omega$-सुसंगत नसलेल्या उपपत्ती आहेत). त्यामुळे
\ref{inc:tcp:oqn:thm:oconsis-q-undec} मधील गृहीतक साध्या सुसंगततेपर्यंत
शिथिल करून अधिक सबल प्रमेय मिळवता येते.
\end{explain}

\begin{lem}
``सार्वत्रिक संगणनक्षम संबंध'' अस्तित्वात नाही. म्हणजे पुढील गुणधर्म
असलेला कोणताही द्विस्थानी संगणनक्षम संबंध $R(x,y)$ नाही:
$S(y)$ हा कोणताही एकस्थानी संगणनक्षम संबंध असेल, तर असा काही $k$
असतो की प्रत्येक $y$ साठी $S(y)$ सत्य असते तर आणि तरच
$R(k,y)$ सत्य असते.
\end{lem}

\begin{proof}
$R(x,y)$ हा सार्वत्रिक संगणनक्षम संबंध आहे असे समजू.
$S(y)$ हा $\lnot R(y,y)$ हा संबंध घेऊ. $S(y)$ संगणनक्षम असल्याने
काही $k$ साठी $S(y)$ हे $R(k,y)$ शी सममूल्य असेल.
पण मग $S(k)$ हे $R(k,k)$ आणि $\lnot R(k,k)$ या
दोन्हींशी सममूल्य ठरेल; हा व्याघात आहे.
\end{proof}


\begin{thm}
$\Th{T}$ ही $\Th{Q}$ समाविष्ट असलेली कोणतीही सुसंगत उपपत्ती
असेल, तर $\Th{T}$ अनिर्णेय असते.
\end{thm}


\begin{proof}
$\Th{T}$ हा $\Th{Q}$ चा सुसंगत, निर्णेय विस्तार आहे असे समजू.
$\Th{T}$ वापरून सार्वत्रिक संगणनक्षम संबंध परिभाषित करू
आणि त्यामुळे व्याघात मिळवू.

$R(x,y)$ पुढील अटीवर, आणि फक्त त्याच अटीवर, सत्य असू द्या:
\begin{quote}
$x$ हा $\theta(u)$ या सूत्राचा संकेतांक आहे, आणि $\Th{T}$ हे
$\theta(\num y)$ सिद्ध करते.
\end{quote}
$\Th{T}$ निर्णेय आहे असे गृहीत धरले असल्यामुळे $R$ संगणनक्षम
आहे: सूत्र-संकेतांक ओळखणे व त्यात संख्यांकाचे प्रतिस्थापन करणेही
संगणनक्षम आहे. आता $R$ सार्वत्रिक असल्याचे दाखवू.
$S(y)$ हा कोणताही संगणनक्षम संबंध असेल, तर त्याचे $\Th{Q}$ मध्ये
(आणि म्हणून $\Th{T}$ मध्येही) $\theta_S(u)$ या सूत्राने निरूपण होते.
म्हणून प्रत्येक $n$ साठी
\begin{eqnarray*}
S(n) & \lif & T \vdash \theta_S(\num n) \\
& \lif & R(\Gn{\theta_S(u)}, n)
\end{eqnarray*}
आणि
\begin{eqnarray*}
\lnot S(n) & \lif & T \vdash \lnot \theta_S(\num n) \\
& \lif & T \not\vdash \theta_S(\num n) \quad \text{($\Th{T}$ सुसंगत असल्यामुळे)} \\
& \lif & \lnot R(\Gn{\theta_S(u)}, n).
\end{eqnarray*}
म्हणजे प्रत्येक $y$ साठी $S(y)$ सत्य असते तर आणि तरच
$R(\Gn{\theta_S(u)}, y)$ सत्य असते. त्यामुळे $R$ सार्वत्रिक
ठरतो आणि अपेक्षित व्याघात मिळतो.
\end{proof}

``सत्य अंकगणित'' ही $\Setabs{\varphi}{ \Sat{\Nat}{\varphi}}$ उपपत्ती
असू द्या; म्हणजे अंकगणिताच्या भाषेत प्रमाण अर्थनिर्धारणानुसार
सत्य असलेल्या वाक्यांचा हा संच आहे.
प्रमाण प्रतिमानात क्यूची स्वयंसिद्धके सत्य असल्याने
सत्य अंकगणित हा त्याचा सुसंगत विस्तार आहे.

\begin{cor}
सत्य अंकगणित अनिर्णेय आहे.
\end{cor}















\section{निर्णेय स्वयंसिद्धकसंचाने निरूपणयोग्य उपपत्ती}\label{inc:tcp:cax:sec}

उपपत्ती $\Th{T}$ साठी स्वयंसिद्धकांचा संगणनक्षम संच~$A$ असेल,
तर तिला \emph{निर्णेय स्वयंसिद्धकसंचाने निरूपणयोग्य} म्हणतात.
($A$ हा $\Th{T}$ चा स्वयंसिद्धकसंच आहे याचा अर्थ
$T = \Setabs{\varphi}{A \Proves \varphi}$.) नैसर्गिक संख्यांचे कोणतेही
``युक्तिसंगत'' स्वयंसिद्धकीकरण हा गुणधर्म पूर्ण करेल.
विशेषतः, सांत स्वयंसिद्धकसंच असलेली कोणतीही उपपत्ती
निर्णेय स्वयंसिद्धकसंचाने निरूपणयोग्य असते.

\begin{lem}
$\Th{T}$ ही निर्णेय स्वयंसिद्धकसंचाने निरूपणयोग्य असेल, तर $\Th{T}$ ही
संगणनक्षमपणे प्रगणनीय असते.
\end{lem}

\begin{proof}
$A$ हा $\Th{T}$ चा संगणनक्षम स्वयंसिद्धकसंच आहे असे समजू.
$\varphi \in T$ असल्याचे होकारार्थी प्रमाण शोधण्यासाठी
स्वयंसिद्धकांपासून $\varphi$ ची निष्पत्ती शोधत राहा;
ते वाक्य सदस्य नसेल, तर हा शोध थांबेलच असे नाही.

थोडे वेगळ्या शब्दांत, $\varphi$ हे $\Th{T}$ मध्ये असते तर आणि तरच $A$ मधील $\psi_1$, \dots, $\psi_k$ अशा स्वयंसिद्धकांची
सांत यादी आणि प्रथम-क्रम तर्कशास्त्रात
$(\psi_1 \land \dots \land \psi_k) \lif \varphi$ ची
निष्पत्ती असते. पण ``असा काही \dots आहे की \dots''
या रूपातील व्याख्या असलेला कोणताही संच संगणनक्षमपणे प्रगणनीय असतो,
जर दुसऱ्या ``\dots'' मधील अट संगणनक्षम असेल, हे आपण आधीच जाणतो.
\end{proof}












\section{निर्णेय स्वयंसिद्धकसंचाने निरूपणयोग्य संपूर्ण उपपत्ती निर्णेय असतात}\label{inc:tcp:cdc:sec}

प्रत्येक वाक्य~$\varphi$ साठी $\varphi$ किंवा $\lnot \varphi$ यांपैकी
किमान एक सिद्ध होत असेल, तर उपपत्तीला {\em संपूर्ण} म्हणतात.

\begin{lem}
उपपत्ती $\Th{T}$ संपूर्ण आणि निर्णेय स्वयंसिद्धकसंचाने निरूपणयोग्य आहे असे समजू.
मग $\Th{T}$ निर्णेय असते.
\end{lem}

\begin{proof}
$\Th{T}$ संपूर्ण आहे आणि $A$ हा तिचा संगणनक्षम स्वयंसिद्धकसंच
आहे असे समजू. $\Th{T}$ विसंगत असेल, तर ती स्पष्टपणे संगणनक्षम
आहे. (वाक्य-संकेतांकासाठी अल्गोरिदम: ``नेहमी होय म्हणा.'')
वाक्याचा संकेतांक नसलेले आदान प्रथम वगळावे. म्हणून $\Th{T}$
सुसंगतही आहे असे मानू शकतो.

वाक्य $\varphi$ हे $\Th{T}$ मध्ये आहे की नाही हे ठरवण्यासाठी,
$\Th{T}$ पासून $\varphi$ ची निष्पत्ती आणि
$\lnot \varphi$ ची निष्पत्ती यांचा एकाच वेळी शोध घ्या.
$\Th{T}$ संपूर्ण असल्यामुळे यांपैकी एक तरी निष्पत्ती मिळेलच;
आणि $\Th{T}$ सुसंगत असल्यामुळे $\lnot \varphi$ ची
निष्पत्ती मिळाल्यास $\varphi$ ची निष्पत्ती असणार नाही.

हेच दुसऱ्या प्रकारे पाहू: $\Th{T}$ ही संगणनक्षमपणे प्रगणनीय आहे, हे आपण
आधीच जाणतो. म्हणून आधी सिद्ध केलेल्या प्रमेयानुसार,
$\Th{T}$ चा पूरक संचदेखील संगणनक्षमपणे प्रगणनीय आहे हे दाखवणे पुरेसे आहे.
पण वाक्य~$\varphi$ हे $\overline{\Th{T}}$ मध्ये असते तर आणि तरच $\lnot \varphi$ हे $\Th{T}$
मध्ये असते. त्यामुळे $\overline{\Th{T}} \leq_m
\Th{T}$.
येथे उपपत्ती म्हणजे वाक्यांच्या संकेतांकांचा संच मानला आहे.
वाक्य-संकेतांकांवर हे न्यूनीकरण निषेध घालून मिळते. सर्व नैसर्गिक
संख्यांवर ते एकूण फलन करण्यासाठी, वाक्याचा संकेतांक नसलेले
प्रत्येक आदान शून्याची स्वतःशी एकरूपता सांगणाऱ्या निश्चित
तार्किक प्रमेयाच्या संकेतांकाकडे पाठवा. अशी आदाने पूरक संचात
असतात आणि त्यांची प्रतिमा उपपत्तीत असते. वाक्य-संकेतांक
ओळखणेही संगणनक्षम असल्याने हे एकूण न्यूनीकरण संगणनक्षम आहे.
\end{proof}












\section[Q चा संपूर्ण, सुसंगत, निर्णेय स्वयंसिद्धकसंचाने निरूपणयोग्य विस्तार नाही]{$\Th{Q}$ चा संपूर्ण, सुसंगत, निर्णेय स्वयंसिद्धकसंचाने निरूपणयोग्य
  विस्तार नाही}\label{inc:tcp:inc:sec}

\begin{thm}
\label{inc:tcp:inc:thm:first-incompleteness}
$\Th{Q}$ चा संपूर्ण, सुसंगत आणि निर्णेय स्वयंसिद्धकसंचाने निरूपणयोग्य विस्तार
अस्तित्वात नाही.
\end{thm}

\begin{proof}
$\Th{Q}$ चा कोणताही सुसंगत, निर्णेय विस्तार नसतो, हे
आपल्याला आधीच माहीत आहे. पण $\Th{T}$ संपूर्ण आणि
निर्णेय स्वयंसिद्धकसंचाने निरूपणयोग्य असेल, तर ती निर्णेय असते.
\end{proof}

\begin{explain}
हे प्रमेय गोडेल यांच्या १९३१ मधील पहिल्या अपूर्णता प्रमेयाच्या
मूळ मांडणीपासून फारसे दूर नाही. अधिक आधुनिक पारिभाषिक शब्द
सोडले, तर मुख्य फरक असे आहेत: गोडेल यांनी ``सुसंगत''ऐवजी
``$\omega$-सुसंगत'' ही अट वापरली; तसेच संगणनक्षमतेची आकारिक
संकल्पना तेव्हा प्रस्थापित झाली नव्हती, म्हणून ``निर्णेय स्वयंसिद्धकसंचाने निरूपणयोग्य''
हे सर्वसाधारणपणे म्हणता येत नव्हते. (पुढील दशकभरात गोडेल
यांच्यासह अनेकांनी संगणनक्षमतेची आकारिक प्रतिमाने विकसित केली;
गोडेल प्रकारची घटना कोणत्या उपपत्तींना लागू होते हे नेमके
सांगणे, हा त्यामागील एक मोठा उद्देश होता.)

संपूर्णता, सुसंगतता आणि निर्णेय स्वयंसिद्धकसंचाने निरूपणयोग्यता हे तिन्ही गुणधर्म
एकाच वेळी मिळू शकत नाहीत, असे प्रमेय सांगते. मात्र यांपैकी
कोणताही एक सोडल्यास उरलेले दोन मिळू शकतात: $\Th{Q}$ सुसंगत
आणि संगणनक्षम स्वयंसिद्धकांनी निरूपित आहे, पण संपूर्ण नाही;
विसंगत उपपत्ती संपूर्ण असून संगणनक्षम स्वयंसिद्धकांनी
(उदाहरणार्थ, $\{ \eq/[0][0] \}$ या संचाने) निरूपित होते,
पण सुसंगत नाही; आणि अंकगणितातील सर्व सत्य वाक्यांचा संच
संपूर्ण व सुसंगत आहे, पण संगणनक्षम स्वयंसिद्धकांनी
निरूपित होत नाही.
\end{explain}











\section[Q मध्ये सिद्ध होणारी आणि ज्यांचे निषेध सिद्ध होतात अशा वाक्यांचे संच संगणनक्षमपणे अविलगनीय आहेत]{$\Th{Q}$ मध्ये सिद्ध होणारी आणि ज्यांचे निषेध सिद्ध होतात
  अशा वाक्यांचे संच संगणनक्षमपणे अविलगनीय आहेत}\label{inc:tcp:ins:sec}


$\overline{\Th{Q}}$ हा ज्या वाक्यांचे {\em निषेध} $\Th{Q}$ मध्ये
सिद्ध होतात त्यांचा संच असू द्या, म्हणजे
$\overline{\Th{Q}} = \Setabs{\varphi}{\Th{Q} \Proves
  \lnot \varphi}$. विभक्त संच $A$ आणि $B$ यांना
संगणनक्षमपणे अविलगनीय म्हणतात, जर असा कोणताही
संगणनक्षम संच $C$ नसेल की $A \subseteq C$ आणि
$B \subseteq \Complement{C}$.

\begin{lem}
$\Th{Q}$ आणि $\overline{\Th{Q}}$ हे संगणनक्षमपणे अविलगनीय आहेत.
\end{lem}

\begin{proof}
$C$ हा असा संगणनक्षम संच आहे असे समजू की
$\Th{Q} \subseteq C$ आणि
$\overline{\Th{Q}} \subseteq \Complement{C}$.
$R(x,y)$ हा पुढील संबंध असू द्या:
\begin{quote}
$x$ हा $\theta(u)$ या सूत्राचा संकेतांक आहे आणि
$\theta(\num y)$ हे $C$ मध्ये आहे.
\end{quote}
सूत्र-संकेतांक ओळखणे, संख्यांकाचे प्रतिस्थापन करणे आणि
या संचातील सदस्यत्व तपासणे संगणनक्षम असल्यामुळे $R(x,y)$
संगणनक्षम आहे. आता तो सार्वत्रिक संगणनक्षम संबंध आहे असे दाखवू;
त्यातून व्याघात मिळेल.

$S(y)$ हा संगणनक्षम संबंध असून त्याचे $\Th{Q}$ मध्ये
$\theta_S(u)$ या सूत्राने निरूपण होते असे समजू. मग
\begin{eqnarray*}
S(n) & \lif & \Th{Q} \Proves \theta_S(\num n) \\
& \lif & \theta_S(\num n) \in C
\end{eqnarray*}
आणि
\begin{eqnarray*}
\lnot S(n) & \lif & \Th{Q} \Proves \lnot \theta_S(\num n) \\
& \lif & \theta_S(\num n) \in \overline{\Th{Q}} \\
& \lif & \theta_S(\num n) \not\in C
\end{eqnarray*}
असे मिळते. म्हणून $S(y)$ हे
$R(\#(\theta_S(u)),y)$ शी सममूल्य आहे.
\end{proof}












\section[Q शी सुसंगत उपपत्ती अनिर्णेय असतात]{$\Th{Q}$ शी सुसंगत उपपत्ती अनिर्णेय असतात}\label{inc:tcp:con:sec}

पुढील प्रमेय असे सांगते की केवळ $\Th{Q}$ अनिर्णेय आहे असे नाही,
तर $\Th{Q}$ शी विसंगती निर्माण न करणारी कोणतीही उपपत्ती अनिर्णेय असते.

\begin{thm}
$\Th{T}$ ही अंकगणिताच्या भाषेतील $\Th{Q}$ शी सुसंगत असलेली
कोणतीही उपपत्ती असू द्या (म्हणजे $\Th{T} \cup \Th{Q}$ सुसंगत आहे).
मग $\Th{T}$ अनिर्णेय असते.
\end{thm}

\begin{proof}
$\Th{Q}$ ची $\mathsf{Q}_1$, \dots, $\mathsf{Q}_8$ ही सांत स्वयंसिद्धके आहेत,
हे आठवा. त्यांच्या जागी $\mathsf{E} = \mathsf{Q}_1 \land
\dots \land \mathsf{Q}_8$ हे एकच स्वयंसिद्धकही घेता येते.

$\Th{T}$ ही $\Th{Q}$ शी सुसंगत असलेली निर्णेय उपपत्ती आहे
असे समजू. पुढील संच घेऊ:
\[C = \Setabs{\varphi}{\Th{T} \Proves \mathsf{E} \lif \varphi}.
\]
निश्चित सूत्राशी अभिव्यंजन तयार करणे संगणनक्षम आहे;
गृहीत उपपत्ती निर्णेय असल्याने या संचातील सदस्यत्वही निर्णेय आहे.
$C$ हा $\Th{Q}$ आणि $\overline{\Th{Q}}$ यांना संगणनक्षमपणे
विलग करणारा संच ठरेल, असे दाखवू; त्यामुळे व्याघात मिळेल.
प्रथम, $\varphi$ हे $\Th{Q}$ मध्ये असेल, तर $\Th{Q}$ च्या
स्वयंसिद्धकांपासून $\varphi$ सिद्ध होते. निगमन प्रमेयानुसार
प्रथम-क्रम तर्कशास्त्रात $\mathsf{E} \lif \varphi$ ची निष्पत्ती
असते. तार्किक प्रमेय असल्याने ते गृहीत निर्णेय उपपत्तीतही सिद्ध होते.
म्हणून $\varphi$ हे $C$ मध्ये आहे.

दुसरीकडे, $\varphi$ हे $\overline{\Th{Q}}$ मध्ये असेल, तर प्रथम-क्रम
तर्कशास्त्रात $\mathsf{E} \lif \lnot \varphi$ ची सिद्धता असते.
$\Th{T}$ हे $\mathsf{E} \lif \varphi$ सुद्धा सिद्ध करत असेल, तर
$\Th{T}$ हे $\lnot \mathsf{E}$ सिद्ध करेल; अशा वेळी
$\Th{T} \cup \Th{Q}$ विसंगत होईल. पण
$\Th{T} \cup \Th{Q}$ सुसंगत आहे असे आपण गृहीत धरले आहे.
म्हणून $\Th{T}$ हे $\mathsf{E} \lif \varphi$ सिद्ध करत नाही आणि $\varphi$
हे $C$ मध्ये नाही.

म्हणजे $\varphi$ हे $\Th{Q}$ मध्ये असेल तर $C$ मध्ये असते,
आणि $\varphi$ हे $\overline{\Th{Q}}$ मध्ये असेल तर $\Complement{C}$
मध्ये असते, हे दाखवले. त्यामुळे $C$ हा संगणनक्षम विलगक
ठरतो; हाच अपेक्षित व्याघात आहे.
\end{proof}

हे प्रमेय अतिशय प्रभावी आहे. उदाहरणार्थ, त्यावरून पुढील निष्कर्ष मिळतो:

\begin{cor}
  अंकगणिताच्या भाषेसाठी प्रथम-क्रम तर्कशास्त्र (म्हणजे
  $\Setabs{\varphi}{\text{$\varphi$ हे प्रथम-क्रम तर्कशास्त्रात सिद्ध होते}}$
  हा संच) अनिर्णेय आहे.
\end{cor}

\begin{proof}
प्रथम-क्रम तर्कशास्त्र हा $\emptyset$ पासून निष्पन्न होणाऱ्या
वाक्यांचा संच आहे; आणि तो रिक्त स्वयंसिद्धकसंच $\Th{Q}$ शी सुसंगत आहे.
\end{proof}












\section[Q च्या अर्थनिर्धारणाशी सुसंगत उपपत्ती अनिर्णेय असतात]{$\Th{Q}$ च्या अर्थनिर्धारणाशी सुसंगत उपपत्ती अनिर्णेय असतात}\label{inc:tcp:itp:sec}

हे निकाल आणखी सबल करता येतात. अनौपचारिकपणे सांगायचे तर,
भाषा $\Lang{L_1}$ चे दुसरी भाषा $\Lang{L_2}$ मध्ये
अर्थनिर्धारण करताना $\Lang{L_1}$ चे विश्व, संबंधचिन्हे आणि
फलनचिन्हे ही $\Lang{L_2}$ मधील सूत्रे वापरून
परिभाषित करावी लागतात. येथे त्यासाठी वेळ देणार नसलो,
तरी या संकल्पनेची नेमकी व्याख्या देता येते.

\begin{thm}
  $\Th{T}$ ही अशा भाषेतील उपपत्ती असू द्या की त्या भाषेत
  अंकगणिताच्या भाषेचे अर्थनिर्धारण करता येते, आणि त्या
  अर्थनिर्धारणातील $\Th{Q}$ शी $\Th{T}$ सुसंगत असते.
  मग $\Th{T}$ अनिर्णेय असते. $\Th{T}$ कडून $\Th{Q}$ च्या
  स्वयंसिद्धकांची अर्थनिर्धारित रूपे सिद्ध होत असतील, तर
  $\Th{T}$ चा कोणताही सुसंगत विस्तार निर्णेय नसतो.
\end{thm}

या प्रमेयाची सिद्धता मागील प्रमेयाच्या सिद्धतेत थोडा बदल करून
मिळते; प्रतिउदाहरणावरून $\Th{Q}$ आणि $\overline{\Th{Q}}$ यांचे
विलगीकरण करता येईल. निवडीच्या स्वयंसिद्धकासह
झर्मेलो--फ्रेंकेल संच उपपत्ती $\Th{ZFC}$ ही सामान्य
गणिताचा मोठा भाग विकसित करण्याइतकी समर्थ स्वयंसिद्धकीय
पायाभूत व्यवस्था मानता येते. $\Th{ZFC}$ मध्ये नैसर्गिक
संख्या परिभाषित करता येतात; आणि या अर्थनिर्धारणातून
$\Th{Q}$ ची स्वयंसिद्धके सत्य ठरतात. म्हणून पुढील निष्कर्ष मिळतो:

\begin{cor}
$\Th{ZFC}$ चा कोणताही सुसंगत, निर्णेय विस्तार नाही.
\end{cor}

\begin{cor}
$\Th{ZFC}$ चा संपूर्ण, सुसंगत आणि संगणनक्षमपणे
निर्णेय स्वयंसिद्धकसंचाने निरूपणयोग्य असा कोणताही विस्तार नाही.
\end{cor}

$\Th{ZFC}$ च्या भाषेत $\in$ हे एकमेव द्विस्थानी संबंधचिन्ह
आहे. (खरे तर समतेचीही गरज नाही.) म्हणून पुढील निष्कर्ष मिळतो:

\begin{cor}
द्विस्थानी संबंधचिन्ह असलेल्या कोणत्याही भाषेसाठी प्रथम-क्रम
तर्कशास्त्र अनिर्णेय आहे.
\end{cor}

दोन एकस्थानी फलनचिन्हे असलेल्या कोणत्याही भाषेलाही हा निकाल
लागू होतो, कारण त्यांचा वापर करून द्विस्थानी संबंधचिन्हाचे
अनुकरण करता येते. हे निकाल नेमक्या सीमेवरचे आहेत: केवळ
\emph{एकस्थानी} संबंधचिन्हे आणि जास्तीत जास्त एक
\emph{एकस्थानी} फलनचिन्ह असलेल्या भाषेसाठी प्रथम-क्रम
तर्कशास्त्र निर्णेय असते, असे सिद्ध होते.

आणखी एक माहितीची बाब. प्रमाण प्रतिमानात सत्य असलेल्या
$\Obj 0$, $'$, $+$, $\times$, $<$ या भाषेतील वाक्यांचा
संच अनिर्णेय आहे, हे आपल्याला माहीत आहे. प्रत्यक्षात,
उरलेल्या चिन्हांच्या साहाय्याने $<$ परिभाषित करता येते,
आणि $\times$ व $'$ यांच्या साहाय्याने $+$ परिभाषित
करता येते. म्हणून $\Obj 0$, $'$, $\times$ या भाषेतील
सत्य वाक्यांचा संच अनिर्णेय आहे. दुसरीकडे, प्रमाण
अंकगणितीय प्रतिमानात सत्य असलेल्या $\Obj 0$, $'$, $+$
या भाषेतील वाक्यांचा संच निर्णेय आहे, हे प्रेस्बुर्गर यांनी
दाखवले. मात्र त्याची निर्णयपद्धती संगणनदृष्ट्या अव्यवहार्य आहे.



